\documentclass{situationist-preprint}

\title{Substantivalism and Newton's Bucket}
\msid{TSM-06-a}
\msfamily{TSM-06}
\msdiscipline{Time, Space and Modality}
\mstype{Principal manuscript}
\mskeywords{substantivalism; relationalism; Newton's bucket; inertial structure; Mach's principle; spacetime; hole argument}

\begin{document}
\maketitle

\begin{abstract}
Newton argued from the concave surface of water in a rotating bucket that rotation is absolute, and hence that there is a space with respect to which bodies move. Relationalists have replied that all motion is relative, that the effect is due to rotation relative to the fixed stars, that absolute acceleration is a primitive property of bodies, or that inertial structure is a regularity in the relations among bodies. We argue that the bucket argument survives these replies when it is stated as an argument for substantivalism about spacetime structure: there is an entity, spacetime, whose affine or metrical structure is not reducible to the relations among material bodies. Each relationalist reply either fails to explain the difference, changes the dynamics, or reintroduces the structure under another name, and in general relativity the structure has degrees of freedom of its own. Space, so understood, is a substance.
\end{abstract}

\section{The problem}

In the Scholium to the Definitions of the \emph{Principia}, Newton describes a bucket of water hung by a twisted cord \citep{Newton1687}. When the bucket begins to spin, the water at first stays flat although it rotates relative to the bucket; as it picks up the rotation it climbs the walls, and its surface is concave when it is at rest relative to the bucket. The concavity is therefore not correlated with rotation relative to the bucket. Newton concluded that it is correlated with \enquote{true and absolute circular motion}, motion relative to absolute space, and he added the case of two globes joined by a cord in otherwise empty space, whose rotation could be detected by the tension in the cord although there is nothing for them to rotate relative to.

\citet{Leibniz1716} held that space is the order of coexisting things and motion the change of relations among bodies. Later relationalists answered the bucket in three ways: the water rotates relative to the fixed stars \citep{Mach1883}; absolute acceleration is a monadic property of bodies \citep{Sklar1974}; inertial structure is a Humean summary of relative motions \citep{Huggett2006}. On the other side, \citet{Stein1967} and \citet{Earman1989} showed that the bucket argues not for Newton's absolute space, with its absolute velocities, but for the affine structure of neo-Newtonian spacetime; \citet{Pooley2013} surveys the replies.

We argue that the bucket argument survives the relationalist replies, and that space, or rather spacetime, is a substance: an entity bearing structure that is not reducible to relations among bodies.

\section{Framework}

\begin{definition}[Substantivalism]\label{d:subst}
\emph{Substantivalism} about spacetime is the thesis that there is an entity, spacetime, with parts (points or regions) that stand in geometrical relations and whose structure, including at least its affine structure, does not supervene on the spatiotemporal relations among material bodies. \emph{Relationalism} is its denial.
\end{definition}

\begin{definition}[Affine structure]\label{d:affine}
The \emph{affine structure} of a spacetime is the structure that distinguishes straight (unaccelerated) worldlines from curved (accelerated) ones; in a neo-Newtonian spacetime it determines which motions are inertial and which bodies rotate.
\end{definition}

\begin{principle}[Dynamical difference]\label{p:dyn}
There are physical effects, such as the concavity of rotating water, the tension in the cord between rotating globes, and the Coriolis deflection, that depend on whether a body moves inertially, and that do not depend on the motion of the body relative to any other body or system of bodies in the cases at issue.
\end{principle}

\begin{principle}[Explanatory realism]\label{p:ibe}
If the best explanation of a pervasive physical difference posits a structure that the difference depends on, and no rival explanation of comparable standing does without it, the structure exists.
\end{principle}

Principle~\ref{p:ibe} is the inference that underwrites scientific realism (SCI-02). The last clause of Principle~\ref{p:dyn} is what the relationalist disputes; Lemma~\ref{l:rel} argues for it.

\section{The argument}

\begin{lemma}[Relative motion does not explain]\label{l:rel}
The effects of Principle~\ref{p:dyn} are not explained by the motion of the rotating body relative to other bodies.
\end{lemma}

\begin{proof}
In the bucket, the effect is absent when relative rotation of water and bucket is greatest and present when it is zero; for the globes, there are no other bodies. More generally, a body at rest with the rest of the universe rotating around it is, in its relations to other bodies, identical to the body rotating and the rest at rest, and the dynamics predicts the effect only in the second case. A purely relational quantity takes the same value in both, so it cannot be what the effect depends on.
\end{proof}

\begin{lemma}[Affine structure explains]\label{l:aff}
Given affine structure, the effects of Principle~\ref{p:dyn} are explained: they occur exactly when a body's worldline is not straight according to that structure, as the dynamical laws state.
\end{lemma}

\begin{proof}
By the law of inertia, force-free bodies follow straight worldlines. The parts of a rotating body follow curved ones and require centripetal forces, supplied in the water by the pressure gradients that make the surface concave and in the globes by the cord.
\end{proof}

\begin{theorem}\label{t:main}
Given Principles~\ref{p:dyn} and~\ref{p:ibe}, and unless a relationalist rival explains the effects with comparable standing, spacetime has an affine structure that does not supervene on relations among bodies; substantivalism is true.
\end{theorem}

\begin{proof}
By Lemma~\ref{l:aff} the affine structure explains the effects, and by Lemma~\ref{l:rel} the relative motions of bodies do not. The structure does not supervene on the relations among bodies, since the two kinematically identical situations in the proof of Lemma~\ref{l:rel} agree in all such relations and differ in affine structure. By Principle~\ref{p:ibe}, the structure exists, and by Definition~\ref{d:subst} substantivalism is true. The proviso is discharged in Section~\ref{s:obj}, where the rival explanations are considered.
\end{proof}

\section{Objections and replies}\label{s:obj}

\begin{objection}[Leibniz]\label{o:leib}
All motion is relative; to say that the whole universe could rotate without any change in the relations among its parts is to posit a difference that makes no difference \citep{Leibniz1716}.
\end{objection}

\begin{reply}
Leibniz's principle applies to differences that make no difference, such as a uniform translation of everything, which neo-Newtonian spacetime does without. Rotation is not one: by Lemma~\ref{l:rel} it makes a dynamical difference. Leibniz conceded that true motion differs from mere change of situation, placing the difference in its cause \citep{Leibniz1716}; what the difference consists in is then not a relation among bodies. TSM-06-b reconsiders his position.
\end{reply}

\begin{objection}[Mach]\label{o:mach}
The water rotates relative to the fixed stars. If the stars were rotated around a still bucket, the water would climb the walls. Newton assumed, without evidence, that only absolute rotation has effects \citep{Mach1883}.
\end{objection}

\begin{reply}
Within Newtonian mechanics the stars rotating around a still bucket would leave the water flat; Mach's proposal calls for a different dynamics, and the candidates confirm the substantivalist reading. General relativity, which Einstein hoped would be Machian, is not: it has vacuum solutions with nontrivial inertial structure and no matter at all, and solutions, such as G\"odel's, in which the matter of the universe rotates relative to the local inertial frames. Frame-dragging is a small correction to an inertial structure the matter does not fix. The relational particle dynamics of \citet{BarbourBertotti1982} is genuinely Machian, but it obtains this by restricting the solutions of Newtonian mechanics to those with zero total angular momentum. That is a law forbidding a rotating universe, not an explanation of inertial effects without inertial structure, and the theory's standing is not comparable to that of general relativity \citep{Pooley2013}.
\end{reply}

\begin{objection}[Sklar's primitive acceleration]\label{o:sklar}
Rotation may be absolute without being motion relative to space: bodies may have a monadic property of absolute acceleration, on which the effects depend \citep{Sklar1974}.
\end{objection}

\begin{reply}
The proposed property must cohere across bodies: the accelerations of a body's parts must agree with its rotation as a whole, and all bodies with zero acceleration must move along one common family of straight lines. On the substantivalist view this follows from the existence of one affine structure; on Sklar's it is an unexplained coincidence among the monadic properties of every body \citep{Earman1989}, and Principle~\ref{p:ibe} prefers the structure.
\end{reply}

\begin{objection}[Regularity relationalism]\label{o:huggett}
Inertial structure is a summary of the regularities in the relative motions of bodies over the whole history of the world, as on the Humean view of laws defended in MET-12-a. It is real as a law is real, but it does not require a substance \citep{Huggett2006}.
\end{objection}

\begin{reply}
The Humean view of laws summarizes a mosaic distributed over a spatiotemporal manifold, which it presupposes; MET-12-a does not reduce geometry to laws. And in general relativity the metric, which determines the affine structure, is a dynamical field with its own degrees of freedom. Gravitational waves propagate in vacuum and carry energy away from binary systems; a structure that does so is not a summary of the bodies' relative motions but something they interact with.
\end{reply}

\begin{objection}[The hole argument]\label{o:hole}
Substantivalism about spacetime points entails that diffeomorphically related models of general relativity represent distinct possibilities, and then the theory is indeterministic in a way that physics does not recognize \citep{EarmanNorton1987}.
\end{objection}

\begin{reply}
The argument tells against a substantivalism on which spacetime points have haecceitistic identities, not against Definition~\ref{d:subst}. The sophisticated substantivalist holds that diffeomorphic models represent the same possibility, because the identity of spacetime points is fixed by their place in the metrical structure \citep{Butterfield1989,Maudlin1988}. Maudlin's metrical essentialism, on which each point's nature includes its metrical relations, fits the thesis of MET-13 that every object has a real definition. The bucket argues for structure, and the hole argument does not touch structure.
\end{reply}

\section{What the result does not establish}

The result does not establish Newton's absolute space: the argument, as corrected by Stein and Earman, establishes affine structure, which in neo-Newtonian spacetime involves no absolute rest and in relativistic spacetime no global simultaneity (TSM-16). It does not establish that spacetime points have identities independent of their structural roles. Nor does it settle whether quantum gravity will treat spacetime as fundamental or emergent.

\section{Conclusion}

Rotation is not a relation among bodies but the curvature of worldlines relative to an affine structure that explains the dynamical difference and does not supervene on such relations. None of the relationalist replies removes the need for it. The bucket argument survives, and space, as spacetime structure, is a substance.

\begin{thebibliography}{99}\small
\bibitem[Barbour and Bertotti(1982)]{BarbourBertotti1982} Barbour, J.~B., and Bertotti, B. (1982). Mach's principle and the structure of dynamical theories. \emph{Proceedings of the Royal Society of London A}, 382, 295--306.
\bibitem[Butterfield(1989)]{Butterfield1989} Butterfield, J. (1989). The hole truth. \emph{The British Journal for the Philosophy of Science}, 40(1), 1--28.
\bibitem[Earman(1989)]{Earman1989} Earman, J. (1989). \emph{World Enough and Space-Time: Absolute versus Relational Theories of Space and Time}. MIT Press.
\bibitem[Earman and Norton(1987)]{EarmanNorton1987} Earman, J., and Norton, J. (1987). What price spacetime substantivalism? The hole story. \emph{The British Journal for the Philosophy of Science}, 38(4), 515--525.
\bibitem[Huggett(2006)]{Huggett2006} Huggett, N. (2006). The regularity account of relational spacetime. \emph{Mind}, 115(457), 41--73.
\bibitem[Leibniz and Clarke(1716)]{Leibniz1716} Leibniz, G.~W., and Clarke, S. (1715--1716). \emph{The Leibniz--Clarke Correspondence}. Ed. H.~G. Alexander, Manchester University Press, 1956.
\bibitem[Mach(1883)]{Mach1883} Mach, E. (1883). \emph{Die Mechanik in ihrer Entwickelung historisch-kritisch dargestellt}. Brockhaus. English translation, \emph{The Science of Mechanics}, Open Court, 1893.
\bibitem[Maudlin(1988)]{Maudlin1988} Maudlin, T. (1988). The essence of space-time. \emph{PSA: Proceedings of the Biennial Meeting of the Philosophy of Science Association}, 1988(2), 82--91.
\bibitem[Newton(1687)]{Newton1687} Newton, I. (1687). \emph{Philosophiae Naturalis Principia Mathematica}. Royal Society. Scholium to the Definitions.
\bibitem[Pooley(2013)]{Pooley2013} Pooley, O. (2013). Substantivalist and relationalist approaches to spacetime. In R.~Batterman (ed.), \emph{The Oxford Handbook of Philosophy of Physics}. Oxford University Press.
\bibitem[Sklar(1974)]{Sklar1974} Sklar, L. (1974). \emph{Space, Time, and Spacetime}. University of California Press.
\bibitem[Stein(1967)]{Stein1967} Stein, H. (1967). Newtonian space-time. \emph{The Texas Quarterly}, 10, 174--200.
\end{thebibliography}

\end{document}
